% 02-reading.tex: как читать карту.
\section{Как читать эту карту}
\label{sec:reading}

Полярис не есть одна огромная таблица или один огромный сервер. Это очень небольшое ядро
удостоверения с намеренно ограниченными структурами, расставленными вокруг него, каждая из
которых стесняет, защищает, наблюдает, проверяет, записывает или ограничивает что-то в структуре,
лежащей внутри неё. Документ читается от ядра наружу, кольцо за кольцом
(рис.~\ref{fig:core-to-shell}), и читатель должен иметь возможность менять увеличение: от
человека к удостоверению, к криптографии, которая его удостоверяет, к правилам схемы, которые его
ограничивают, к органу, который его выдал, к сторонам, которые на него полагаются, к федерации
органов, к свидетелям, которые за ними наблюдают, к межведомственному обмену, к операциям,
которые поддерживают в нём жизнь, к обществу, внутри которого он находится, к среде программных
агентов, которую он теперь обслуживает, и наконец наружу: к программам, написанным посторонними,
чьё принятие или отказ есть единственное свидетельство, которое репозиторий произвёл не о себе
сам.

\subsection{Вопросы, задаваемые каждому слою}

О каждой крупной составной части документ отвечает на одни и те же вопросы, в карточке того вида,
который используется повсюду. Что это такое? Почему оно существует? Какую задачу оно решает? Чему
оно доверяет? Чему оно доверять отказывается? Что оно видит? Чего ему намеренно не дают видеть?
Что происходит, если оно отказывает? Какая часть Полярис его проверяет? Нужно ли оно сейчас или
это подготовка к будущей среде? Как оно соединено со слоем внутри него и со слоем снаружи? Читатель,
способный ответить на эти вопросы про каждое кольцо, владеет той мысленной моделью, ради построения
которой настоящий документ и существует.

\subsection{Город}

\analogy{Запись удостоверения есть реестр, хранимый в городской ратуше: по одной записи на
человека, утверждение, вокруг которого устроено всё остальное. Стена есть схема: триггеры,
ограничения целостности и уникальные индексы, решающие, какие состояния вообще вправе
существовать, и связывающие одинаково каждого посетителя, каждого чиновника и каждое
восстановление из резервной копии. Ворота есть конечные точки, через которые утверждение входит и
выходит: выдача, проверка, вход, обмен. Сторожевые башни есть журналы прозрачности, а свидетели
стоят за стеной, совместно подписывая то, что башни публикуют, и поднимая крик, когда две башни
расходятся. Другие города есть другие удостоверяющие органы, у каждого своя ратуша и свой ключ, и
соединены они дорогами, которые идут в одну сторону и ради одной цели. Надо всем стоят законы,
которых город не в силах отменить голосованием своих чиновников. А за стенами есть путники из
городов, которых здесь никто не строил и которые ничем не обязаны этому городу; то, что они найдут у
его ворот, есть последнее, что показывает эта карта.}

Карта города, с техническим именем рядом с каждым городским, есть рис.~\ref{fig:town-map}
ниже. Аналогия будет возвращаться в затенённой врезке вроде помещённой выше, и
за каждым возвращением следует механизм, который она обозначает. Каждое кольцо к тому же
открывается одной строкой курсивом: единственной мыслью документа, заново сказанной для этого кольца,
о том, что кольцо внутри него было верным и всё же недостаточным.

\figfit{0.55}{town-map}{Обнесённый стеной город, с настоящей составной частью рядом с каждым городским
именем. Реестр в ратуше есть ядро удостоверения; стена есть схема; четверо ворот суть маршруты,
через которые утверждение входит и выходит; сторожевые башни суть журналы прозрачности, а
свидетели за стеной за ними наблюдают; дорога во второй город есть направленная аттестация,
привязанная к контексту; законы над городом связывают сам город. План описывает устройство.
Полярис федеративен: каждый орган есть свой собственный город.}{fig:town-map}


\figfile{core-to-shell}{Концентрическая карта Полярис. Ядро есть одна запись удостоверения и её
подпись. Каждое кольцо называет действительные составные части на нём, а порядок чтения документа
следует кольцам наружу. Внешнее кольцо показано пунктиром, потому что по большей части оно не
код: оно есть то, что стороны за пределами репозитория сделали с кольцами внутри него, а на
\code{1.0.0-кв.7} это один кошелёк, один размещённый прогон проверки соответствия и один
опубликованный набор тестов.}{fig:core-to-shell}

\subsection{Доверие и отказ}

Самая полезная привычка при чтении Полярис состоит в том, чтобы спрашивать о каждом слое, что
именно он отказывается принимать на веру от слоя под ним. Схема отказывается доверять приложению.
Отделяемый верификатор отказывается доверять серверу эмитента. Второй свидетель отказывается
доверять первой реализации. Доверяющая сторона отказывается доверять транзитивной цепочке
удостоверяющих органов. Свидетель отказывается доверять собственной вершине журнала. Слой
инвариантных проверок отказывается доверять проверке, которая не способна обнаружить собственное
нарушение. Мутационные учения отказываются доверять тесту, которому ни разу не показали, как он
падает. А сводная таблица отказывается засчитывать что бы то ни было, что автор прогнал против
собственного кода. Заключение рисует эти отказы лестницей (рис.~\ref{fig:layer-ladder});
разделы~\ref{sec:core} по \ref{sec:operations} объясняют каждый из них.
